\section{The exact global comparison lattice}
\label{sec:comparison-span}

\status{CS1--CS9 prove the integral descent, local comparison and
global span arguments for the supplied compatible forms. In particular,
the source's \texttt{ValIntRec} and comparison-lattice equality are
deduced below; they are not new assumptions. The existence of the
polynomial values \texttt{ValEx}, and the simultaneous forms and carrier
data on which those values depend, remain stated inputs.}

\paragraph{Setting and notation.}
Keep the CM field $E/F$, the same global frame $B$ from SF, and the
same embeddings and positions from LI. Write $L$ for the coefficient
number field, $e:L\hookrightarrow K$ for the prescribed-place embedding,
$O\subset K$ for the local DVR and
$R=L\cap e^{-1}(O)=\mathcal O_{L,w}$.
Let $V$ be the global space of level-zero forms and
$\mathcal L\subset V$ its specified full integer lattice. The comparison
map $C:V\to\mathcal F_K$ is semilinear for $e$; its values are polynomial
functions of the local lattice tuples. The target $M\subset\mathcal F_K$
is the literal module \texttt{G.M}: its values lie in the integral
\emph{Weyl} tensor lattices and satisfy the prescribed level-character
equivariance. The Weyl lattice $\Delta$ is generated by divided-power
words on the highest polynomial. It is not replaced by the coefficient
lattice $\nabla$ of IW.

The intended conclusion is an equality of integral modules:
\[
 \boxed{\quad O\text{-span}\,C(\mathcal L)=M.\quad}       \tag{CS.1}
\]
Equality after taking the $K$-span is only one step in its proof.

\paragraph{CS1: exact descent through a valuation ring.}
Let $L\subset K$ be any field extension, $O$ a valuation ring of $K$,
and $R=O\cap L$. For polynomials $q_i,Q$ over $L$,
\[
 Q\in O\text{-span}\{q_i\}\quad\Longleftrightarrow\quad
 Q\in R\text{-span}\{q_i\}.                            \tag{CS.2}
\]
Only finitely many $q_i$ occur in a given expression. If they are
linearly independent over $L$, an invertible coefficient minor shows
that their coefficients in an expression for $Q$ lie in $L$; hence
they lie in $R$. Otherwise choose a nonzero relation
$\sum_i a_iq_i=0$. The valuation-ring property permits a pivot $j$
with $a_j\ne0$ and $a_i/a_j\in R$ for every $i$.
In $Q=\sum_i c_iq_i$, replace
\[
 c_i\longmapsto c_i-c_j a_i/a_j.
\]
All new coefficients remain in $O$ and the $j$th becomes zero.
Induction proves (CS.2); the reverse inclusion is immediate.
This is the source's \texttt{descVal}. It does not require $K/L$ to
be finite or $O$ to be a finite $R$-module.

\paragraph{CS2: the integral Weyl comparison, including odd positions.}
For a decreasing integral weight $\lambda=(\lambda_1,\ldots,\lambda_n)$,
put $\nu_i=\lambda_i-\lambda_n$ and
$\lambda^*=(-\lambda_n,\ldots,-\lambda_1)$. Let $R_n$ be the
column-reversal matrix and $\epsilon=\det R_n\in\{1,-1\}$.
On invertible matrices define
\[
 (\mathcal D_\lambda f)(X)=
 \epsilon^{-\lambda_1}\det(X)^{\lambda_1-\lambda_n}
                      f(X^{-\mathsf T}R_n).             \tag{CS.3}
\]
It restricts to an integral isomorphism
$\Delta_{\nu(\lambda^*)}(\mathbb Z)\simeq
\Delta_{\nu(\lambda)}(\mathbb Z)$.
Here is a direct proof. Jacobi's complementary-minor identity sends
each initial minor of $X^{-\mathsf T}R_n$ to a signed complementary
minor of $X$ divided by $\det X$. In the highest polynomial the total
denominator exponent is $\lambda_1-\lambda_n$, so (CS.3) sends it
to a signed row-reversal of the highest polynomial for $\lambda$.
For $a\ne b$, substitution by $I+tE_{ab}$ intertwines with substitution
by $I-tE_{ba}$. Comparing powers of the indeterminate $t$ intertwines
their divided powers integrally, with signs. Row permutations preserve
the lattice: signed transpositions are products of integral elementary
matrices, and diagonal sign changes act integrally on every weight.
Diagonal divided powers act by integer binomial coefficients on weight
vectors. Thus all generating words have polynomial integral images.
Applying (CS.3) twice, with $\lambda$ and $\lambda^*$, is the identity;
this proves both inclusions, including all signs.

Equivalently, if $\widetilde f=\mathcal D_\lambda f$, then
\[
 \det(Z)^{\lambda_n}\widetilde f(Z)=
 \det(Z^{-\mathsf T}R_n)^{-\lambda_1}
                                    f(Z^{-\mathsf T}R_n).       \tag{CS.4}
\]
This is the literal \texttt{OddId}/\texttt{OddFact} identity. Negative
determinant twists are allowed. They are evaluated only on invertible
matrices; no determinant denominator is inserted into $\Delta$.

At a lattice $A O^n$, the integral factor used here is exactly
\[
 D_\lambda(AO^n)=
 \det(A)^{-\lambda_n}\,R(A^{-1})\Delta_\nu(O),\qquad
 (R(g)f)(X)=f(gX).                                      \tag{CS.5}
\]
The divided-power generators have integer coefficients and commute
with coefficient extension. Hence (CS.5) and every tensor product of
these factors are the $O$-spans of their corresponding integer
generators after the displayed substitutions. In even positions use
the basis matrix $N_\tau$; in odd positions use $N_\tau^{-\mathsf T}$
and (CS.4), with the source's exact position reindexing: $\lambda_\tau=\lambda^0_{v,s}$,
while the local factor has weight $\lambda^0_{v,s}$ for even $s$ and
$(\lambda^0_{v,s})^*$ for odd $s$. This is \texttt{WtTie}. CS1 descends their coefficients to $R$.
This proves \texttt{WeylLatticeDescent}: membership of a coefficient-field
polynomial in the local Weyl tensor lattice is equivalent to an
$R$-linear combination of the global integer Weyl tensors in the
corresponding frame. Equality of the functions on all invertible complex
frames implies equality of the polynomials: substitute $X+tI$ in each
block and use that only finitely many $t$ give a singular block.

\paragraph{CS3: the odd-position lattice has no hidden index.}
Fix a class lattice $L_c$, its adelic representative $y_c$, and a
matrix $B'$ which is a basis of $L_c^\#$ at all primes over $p$.
Such a matrix exists. The ring obtained from $\mathcal O_E$ by
inverting integers prime to $p$ is a semilocal Dedekind domain;
each of its ideals is principal, by prescribing the finitely many
valuations using CRT. Its finite torsion-free modules are free.
Applied to the full Hermitian dual lattice, this gives exactly
\texttt{IsLocBasis}, whose denominators are integers prime to $p$.

At $v\mid p$, $p\nmid T$ and the source inclusions
$L_b\subset L_b^\#$, $T L_b^\#\subset L_b$ show that $L_b$ is
locally self-dual. The unitary $y_c$ carries it to $L_c$, so the same
holds for $L_c$ and for the lattice spanned locally by $B'$.
Consequently its Gram matrix
\[
 G'=B'^*JB',\qquad X^*=\bar X^{\mathsf T},
\]
and $G'^{-1}$ have integral entries at these places. The first assertion
tests pairs of basis vectors; the second tests the dual basis and uses
self-duality. The \texttt{Joint0} identity identifies the level lattice
with the one framed by $B$, so this argument concerns the same local
identification as SF6.

For a comparison point $\gamma y_c\kappa$ with $\gamma^*J\gamma=J$
and $\kappa$ in the level group, put
\[
 N=B^{-1}\gamma B'.
\]
The point identity and $\kappa$'s preservation of the base lattice give
$L_v=\operatorname{colspan}_{W_v}i_v(N)$.
Indeed the change-of-basis matrix between the two descriptions and its
inverse preserve the integral coordinate module; both are integral,
so their column spans agree. Now $B^*JB=I$ gives
\[
 N^*N=G',\qquad (N^{-1})^*=NG'^{-1}.                    \tag{CS.6}
\]
For an even position, LI's compatibility identifies the two matrices
entry by entry. For an odd position it identifies the coefficient
matrix $e(t_\tau(N)^{-\mathsf T})$ with the Frobenius image of
$(N^{-1})^*$. By (CS.6) that differs from the Frobenius image of $N$
by right multiplication by the integral invertible matrix $G'^{-1}$.
Thus its $O$-column span is identical. This proves the literal
\texttt{RestLatW}, \texttt{RestGramW} and \texttt{RestLatO} clauses.
The $c^s$ in LI is essential to (CS.6); deleting it would give the
wrong odd-position matrix.

\paragraph{CS4: comparison preserves exactly the specified integrality.}
Supply the source's polynomial-value condition \texttt{ValEx}.
It says that the global form at a comparison point is represented by
one coefficient-field polynomial $Q$, after the specified frame
substitution and level character. CS2 and CS3 identify the membership
condition for $e(Q)$ in the local Weyl tensor lattice with the global
integer-Weyl expression in the basis $B'$.

The character factor does not change this condition. Under
\texttt{LevelZeroCharacterRational}, its factors have representatives
$r\in L$ with $r^a=1$ for a positive integer $a$: they come from
additive characters of finite residue fields. Hence
$a\,v_w(r)=0$, so both $r$ and $r^{-1}$ belong to $R$.
Multiplication by the character, or its inverse, preserves the
$R$-span of integer Weyl tensors. Translating the frame by
$\gamma$ and then by $\gamma^{-1}$ proves both directions of
\texttt{valIntAt\_of}; using one local basis in the reverse direction
and every local basis in the forward direction proves its exact
quantifiers. In the actual record this is \texttt{ValIntRec}.

\paragraph{CS5: the localization criterion is an equality.}
The global lattice $\mathcal L$ is defined by the specified Weyl
integrality conditions at every finite coefficient place. Suppose
$x\in V$ satisfies them at $w$. Choose $0\ne d\in\mathcal O_L$
with $dx\in\mathcal L$, using fullness and finite generation.
By CRT choose $s\in\mathcal O_L$ with
\[
 s\equiv1\pmod w,\qquad
 v_{w'}(s)\ge v_{w'}(d)\quad(w'\ne w).                 \tag{CS.7}
\]
Only primes dividing $d$ impose nontrivial conditions. At $w$, $sx$
is integral because $s$ is; elsewhere $sx=(s/d)(dx)$ is integral.
Thus $sx\in\mathcal L$, while $s^{-1}\in R$. This proves
$x\in R\mathcal L$. The converse follows by linearity of every
local condition. CS4 and the comparison's point/class surjectivity
therefore give
\[
 \boxed{\ C(x)\in M\quad\Longleftrightarrow\quad
                         x\in R\mathcal L.\ }          \tag{CS.8}
\]
The level-equivariance part follows from the same point identity and
the matching inverse away character from LI8. Polynomial uniqueness
makes the comparison independent of a chosen point representative.

\paragraph{CS6: one global matrix realizes every local lattice tuple.}
For the finitely many distinct places $v\mid p$, take arbitrary
$g_v\in\mathrm{GL}_n(F_v)$. Choose $a$ so large that
$p^a g_v^{-1}$ is integral for every $v$. Simultaneous approximation
gives $A\in\operatorname{Mat}_n(F)$ with
\[
 A_v-g_v\in p^{a+1}\operatorname{Mat}_n(\mathcal O_{F_v}).
\]
For completeness, approximate each entry separately in its completion,
clear their finitely many global denominators with one integer, and
apply CRT to the distinct prime-power ideals with correspondingly
increased precision. Dividing back gives the displayed approximation.
Now $g_v^{-1}A_v=I+pD_v$, which is invertible over
$\mathcal O_{F_v}$ by reduction modulo $p$. Hence
\[
 A_v\mathcal O_{F_v}^{n}=g_v\mathcal O_{F_v}^{n}.    \tag{CS.9}
\]
Its determinant is nonzero locally, so $A\in\mathrm{GL}_n(F)$.
The empty tuple permits $A=I$. Via LI's completion maps this proves
\texttt{WeakApprox} and \texttt{RatTup}: the approximation produces
\emph{exact} lattice equality, not an approximation to the desired
integral lattice.

\paragraph{CS7: rational integral functions span the local module.}
Call a local polynomial or function rational if every coefficient
belongs to $e(L)$. By CS6 the tuple lattices can be represented by one
global matrix. LI's positions show that all its Frobenius-twisted
matrix entries belong to $e(L)$. The integer divided-power generators
in CS2 therefore show that every local Weyl tensor lattice is the
$O$-span of rational elements.

This also holds for its stabilizer-invariant part. Here is a proof
that does not divide by a stabilizer order. Choose finitely many
rational generators and let $D_R$ be their $R$-span in the finite
$L$-space $U_L$. It is finite free. Let $U_0\subset U_L$ be the
subspace cut out by the rational level-character equations. These
equations are rational because the matrices are global and the
characters have the coefficient-field representatives from LI8.
The intersection $D_R\cap U_0$ is saturated in $D_R$, so it is a
direct summand over the DVR $R$. Clearing a common $R$-denominator
shows that its $L$-span is exactly $U_0\cap LD_R$. Choose a basis
adapted to this summand. The same basis over $O$ proves
\[
 (O D_R)\cap(K U_0)=O(D_R\cap U_0).                    \tag{CS.10}
\]
The common kernel of even infinitely many rational equations commutes
with this field extension: their rows span a finite-dimensional space,
so a finite subset suffices. Thus (CS.10) is exactly the invariant
claim, without an averaging denominator.

There are finitely many component/lattice orbits. Evaluate a function
at one representative per orbit and extend each value equivariantly,
setting it zero on the other orbits. A stabilizer-invariant value makes
this well defined. The sum of these finitely many extensions recovers
the original function. Rational values extend to rational functions
because the character-twisted action is rational. It follows that
\[
 M=O\text{-span}\{F\in M:F\text{ is rational}\}.        \tag{CS.11}
\]
This is the source's \texttt{MRatSpan}.

\paragraph{CS8: rank equality gives rational surjectivity.}
The comparison is injective: every global class has a comparison point,
and if its polynomial value vanishes, the value identity forces the
global form to vanish on every invertible complex frame. The global
realization of forms is itself injective. Its values have rational
coefficients by construction.

An $L$-independent family in $V$ stays $K$-independent after comparison.
Indeed, expand the coefficients of a finite $K$-relation in an
$L$-basis of their span. Taking each basis coordinate in every
polynomial coefficient produces an $L$-relation. Injectivity kills
each relation separately. The supplied rational-forms record identifies
the complex images of a basis of $V$ with a basis of \texttt{Fi(0)}
(its \texttt{emb\_basis} clause), and \texttt{Fi\_zero} identifies
that space with $\mathcal M_{u,0}$. Thus $\dim_LV=\dim_{\mathbb C}\mathcal M_{u,0}$.
CR6 now gives $\dim_LV=\operatorname{rank}_O M$; this is the complete
\texttt{card\_eq\_d} bridge. CS5 shows the basis images
belong to $KM$, because $\mathcal L$ spans $V$. Finiteness of $M$
then shows these images form a $K$-basis of $KM$.

If $F\in M$ is rational, write $F=\sum_i a_iC(b_i)$ with $a_i\in K$.
Choose an $L$-linear retraction $\pi:K\to L$ with $\pi(1)=1$,
and apply it coefficient by coefficient. Rationality gives
\[
 F=C\Bigl(\sum_i\pi(a_i)b_i\Bigr).                     \tag{CS.12}
\]
This proves \texttt{surjRat\_pf}. It is stronger than rational-span
equality in the precise way required next, but it does not yet assert
integral surjectivity.

\paragraph{CS9: the integral span and its exact hypotheses.}
For $x\in\mathcal L$, CS5 gives $C(x)\in M$, proving the easy
inclusion in (CS.1). Conversely, by CS7 it suffices to consider a
rational $F\in M$. CS8 writes $F=C(x)$, and CS5 gives
$x\in R\mathcal L$. Since $e(R)\subset O$, semilinearity implies
$C(x)\in O\text{-span}\,C(\mathcal L)$. This proves (CS.1).

For the literal source data, the theorem just proved is
\texttt{comparison\_lattice\_span\_eq} from
\texttt{CmpCore}, \texttt{GatIs} and \texttt{ValIntRec}.
CS1--CS4 supply the last premise from the same frames and positions,
supplied \texttt{Forms0} and \texttt{Joint0}, coefficient-field character
rationality, and \texttt{CmpCore.ex}.
\texttt{GatIs} identifies the target with the actual local group
model, supplying its finite class set, fraction field and finite
Weyl lattice. CR supplies the exact rank equality. Thus the remaining
comparison input is the compatible global forms/value construction,
including \texttt{ValEx} and the supply of \texttt{Joint0}, not a second
unproved integral-span assertion. For supplied \texttt{Joint0},
\texttt{SameUloc} is its literal local-subgroup equality and needs
no further theorem.

Nothing here identifies Hermitian norms or covolumes. The exact
integral module equality is a necessary input to the later metric
transport and ramified restrictions; those metric and arithmetic
window statements retain their separate status.
